\documentclass[11pt]{article}
\usepackage[a4paper,margin=28mm]{geometry}
\usepackage{amsmath,amssymb,amsthm,mathtools,hyperref}
\newcommand{\ddc}{dd^c}
\newtheorem{theorem}{Theorem}
\newtheorem{lemma}[theorem]{Lemma}
\theoremstyle{remark}\newtheorem{remark}[theorem]{Remark}
\title{Positive fixed-density volume at a nonbig\newline Chern--Ricci flow endpoint}
\author{}
\date{}
\begin{document}
\maketitle
\begin{abstract}
We exhibit a finite-time Chern--Ricci flow whose nef limiting background
has strictly positive fixed-density Gauduchon volume and is not big.
On $\mathbb P^{n-1}\times\mathbb P^1$, $n\geq3$, we construct a smooth
semipositive nonclosed form with exact variational volume $n$, although
its vertical coefficient vanishes on one fiber. We realize its smooth
gauge class as the actual endpoint background of a smooth positive
Hermitian initial metric. The maximal-time theorem of Tosatti--Weinkove
identifies the endpoint time, including the normalization change.
An explicit normalized closed current on the exceptional fiber
annihilates the endpoint and has zero determinant measure. We also specify
the support and matrix of every annihilator in this example and the exact
averaged trace identity at arbitrary finite flow endpoints. No local
convergence theorem on a big locus is contradicted or claimed.
\end{abstract}

We use $\ddc=\frac{i}{2\pi}\partial\bar\partial$, ordinary wedge powers,
and genuine canonical quasi-psh representatives. For a smooth positive
probability $\mu$ and every smooth positive Gauduchon metric $\gamma$, set
\[
 D_\mu(\gamma)=\int_X(\mu/\gamma^n)^{1/n}\gamma^n,
 \quad \lambda_\mu(\theta)=\inf_{\ddc\gamma^{n-1}=0}
 \frac{\int_X\theta\wedge\gamma^{n-1}}{D_\mu(\gamma)},
 \quad V_\mu(\theta)=\lambda_\mu(\theta)^n.
\]
All such metrics participate. At a nef boundary the equivalent decreasing
positive regularization is understood. A smooth background is nef if it
admits arbitrarily small smooth positive perturbations by $\ddc$-potentials,
and big if $\theta+\ddc v\geq\delta h$ for a global quasi-psh $v$, a
Hermitian $h$, and $\delta>0$. Neither definition assumes closedness.
Our flow convention is $\partial_t\omega=-\operatorname{Ric}(\omega)$,
where $\operatorname{Ric}(\omega)=-\ddc\log\omega^n$.

\begin{theorem}\label{th:flow}
For every $n\geq3$ the construction below gives a smooth positive
Hermitian initial metric $\omega_0$ on $X=\mathbb P^{n-1}\times\mathbb P^1$
whose smooth Chern--Ricci flow has finite maximal time $\tau=1$.
For a smooth strictly positive probability $\mu$, its actual limiting
background
\[
 \theta_\tau=\omega_0-\operatorname{Ric}(\omega_0)
\]
is nef and nonbig, while $V_\mu(\theta_\tau)=n>0$.
There is a positive closed $(n-1,n-1)$ current $T$ with
$\int_X\omega_0\wedge T=1$, $\int_X\theta_\tau\wedge T=0$, and
zero determinant measure. This is a counterexample to the
unrestricted numerical-to-big implication for actual flow endpoints.
\end{theorem}

Put $m=n-1\geq2$, $\beta=(m-1)/m$, $B=\mathbb P^m$ and
$X=B\times\mathbb P^1$. Normalize Fubini--Study forms by
$\int_B\omega^m=\int_{\mathbb P^1}w=1$, and put
\[
 x([Z])=\frac{\sum_{j=1}^m|Z_j|^2}{\sum_{j=0}^m|Z_j|^2},
 \quad a=x/\beta,\quad h=\omega+w,\quad\mu=\omega^m\wedge w,
 \quad\theta=\omega+a w.
\]
For positive $\epsilon$, write
$H_\epsilon=m\int_0^1y^{m-1}/(y/\beta+\epsilon)\,dy$.
Set $\theta_\epsilon=\theta+\epsilon w=\omega+(a+\epsilon)w$.
All pullbacks are implicit. The following construction is explicit and
self-contained; no companion result is a trusted premise.

\section{Global smooth radial positive forms}
The function $x$ is smooth globally, ranges in $[0,1]$, vanishes only at
$p=[1:0:\cdots:0]$, and equals one on the hyperplane $Z_0=0$.
In the chart $Z_0=1$, write $s=|z|^2$; then $x=s/(1+s)$ and
$\omega=\ddc\log(1+s)$.
Polar integration of the Fubini--Study determinant $(1+s)^{-m-1}$,
with normalization $\int\omega^m=1$, gives
\begin{equation}\label{eq:radial}
 \int_B b(x)\omega^m=m\int_0^1 b(y)y^{m-1}\,dy.
\end{equation}
Indeed the radial density in $s$ is
$m s^{m-1}(1+s)^{-m-1}ds$, which becomes $m x^{m-1}dx$.
In particular $\int_B\omega^m/a=m\beta/(m-1)=1$.

For $\epsilon>0$ use the preceding $H_\epsilon$ and define
\begin{equation}\label{eq:flux}
 q_\epsilon(x)^m=\frac{m}{H_\epsilon}
 \int_0^x\frac{y^{m-1}}{y/\beta+\epsilon}\,dy,
 \qquad
 \varphi_\epsilon(x)=\int_0^x
 \frac{q_\epsilon(y)-y}{y(1-y)}\,dy.
\end{equation}
The positive real root is used. At zero,
\[
 q_\epsilon(x)=x\left(\frac{m}{H_\epsilon}
 \int_0^1\frac{u^{m-1}}{xu/\beta+\epsilon}\,du\right)^{1/m},
\]
so it is $x$ times a smooth positive function. At one,
$q_\epsilon(1)=1$. Both apparent singularities in
$\varphi_\epsilon'$ are removable smoothly. Thus
$\varphi_\epsilon(x)$ is a global smooth base function.

For any radial local potential $F(s)$, put $q=sF'(s)$.
The tangential Hessian eigenvalue is $F'=q/s$ and the radial one is $dq/ds$.
For $F_\epsilon=\log(1+s)+\varphi_\epsilon(s/(1+s))$, direct
differentiation of \eqref{eq:flux} gives $sF_\epsilon'=q_\epsilon(x)$.
Both eigenvalues are positive, including at $s=0$ by the preceding smooth
factorization. If $A_\epsilon=\omega+\ddc\varphi_\epsilon$, the determinant
ratio in the affine chart is
\begin{equation}\label{eq:density}
 \frac{A_\epsilon^m}{\omega^m}
 =\frac{q_\epsilon^{m-1}\,dq_\epsilon/dx}{x^{m-1}}
 =\frac1{H_\epsilon(a+\epsilon)}.
\end{equation}
At the hyperplane $x=1$ the globally smooth form is the continuous limit
of positive forms, and its determinant ratio in \eqref{eq:density} tends
to a strictly positive number. Thus it is strictly positive there too.
Equivalently its tangential coefficient is that of $\omega$, and its normal
coefficient is $(dq_\epsilon/dx)(1)=1/[H_\epsilon(1/\beta+\epsilon)]>0$.
This proves global smooth strict positivity. It is closed by construction.

\section{A genuine bounded limiting potential}
By monotone convergence $H_\epsilon\to1$ as $\epsilon\downarrow0$.
For each $0<x\leq1$, \eqref{eq:flux} then gives
$q_\epsilon(x)\to x^\beta$.
We justify the limit of the potentials at the singular point as well.
Relative to the probability density $m y^{m-1}dy$, the weight
$(y/\beta+\epsilon)^{-1}$ is decreasing. Its normalized cumulative mass
is therefore at least the unweighted cumulative mass $x^m$.
Relative to the endpoint probability density $(m-1)y^{m-2}dy$, its weight
ratio is proportional to $y/(y+\beta\epsilon)$, which is increasing; its
normalized cumulative mass is at most $x^{m-1}$.
These two elementary comparisons follow by splitting either probability
at $x$: a decreasing weight has at least as large an average on $[0,x]$
as on $[x,1]$, and an increasing weight has the reverse inequality.
Taking positive $m$th roots yields
\begin{equation}\label{eq:bounds}
 x\leq q_\epsilon(x)\leq x^\beta.
\end{equation}

Define
\[
 \varphi(x)=\int_0^x\frac{y^\beta-y}{y(1-y)}\,dy.
\]
Its derivative is integrable at zero because $\beta>0$, and extends smoothly
at one with value $1-\beta$. Hence $\varphi$ is globally bounded continuous,
smooth off $p$, and $\varphi(p)=0$.
Equation~\eqref{eq:bounds} supplies the integrable derivative majorant
$(y^\beta-y)/[y(1-y)]$. Dominated convergence in $L^1([0,1])$ gives
\[
 \sup_{[0,1]}|\varphi_\epsilon-\varphi|\longrightarrow0.
\]
Consequently $A=\omega+\ddc\varphi$ is a positive closed current by
distributional convergence. Local smooth potentials for $\omega$ plus
$\varphi$ are continuous uniform limits of psh functions, so they are
genuinely psh, with the canonical value at $p$. This establishes admissibility
on the whole compact base, not only on its punctured affine chart.

Off $p$ the limiting radial flux is $q=x^\beta$, and the determinant formula
gives
\begin{equation}\label{eq:enddensity}
 \frac{A^m}{\omega^m}=
 \frac{(x^\beta)^{m-1}\beta x^{\beta-1}}{x^{m-1}}=\beta/x=1/a.
\end{equation}
The form is smooth strictly positive there, including at infinity by
the same determinant and continuity argument as before.
No power of this current at $p$ is needed in the following all-test bound.

\section{Exact volume over every smooth Gauduchon test}
On $X$, the probability $\mu$ is smooth strictly positive and
$\theta=\omega+a w$ is smooth semipositive. Thus $\theta+\eta h>0$
for every $\eta>0$, with zero perturbation potential, proving nefness.
Also $d\theta=da\wedge w\ne0$.
The bounded canonical potential pulled back from $B$ gives the positive
current $\Omega=\theta+\ddc\varphi=A+a w$. Off
$Z=\{p\}\times\mathbb P^1$ it is smooth strictly positive, and
\[
 \Omega^n=n a A^m\wedge w=n\mu.
\]
For every smooth positive Gauduchon $\gamma$, put $S=\gamma^{n-1}$.
Distributional Stokes gives $\int_X\theta\wedge S=\int_X\Omega\wedge S$,
because $\ddc S=0$ and $\varphi$ is bounded.
Matrix AM--GM off $Z$ and positivity of the global left measure imply
\begin{equation}\label{eq:lower}
 \int_X\theta\wedge\gamma^{n-1}\geq
 n^{1/n}\int_X(\mu/\gamma^n)^{1/n}\gamma^n.
\end{equation}
The set $Z$ has zero smooth volume; any singular left measure there only
strengthens this inequality. Thus $\lambda_\mu(\theta)\geq n^{1/n}$
over the entire original family. Using the closed metric $h$ shows finiteness.

We next give exact smooth minimizing tests for the positive perturbations.
Set $c_\epsilon=n/H_\epsilon$ and
\[
 \Omega_\epsilon=A_\epsilon+(a+\epsilon)w,
 \qquad\gamma_\epsilon=(a+\epsilon)^{-1/m}\Omega_\epsilon.
\]
These are smooth strictly positive. Equation~\eqref{eq:density} gives
$\Omega_\epsilon^n=c_\epsilon\mu$, and
\[
 T_\epsilon=\gamma_\epsilon^m
 =(a+\epsilon)^{-1}A_\epsilon^m
       +m A_\epsilon^{m-1}\wedge w
\]
is closed: its first term is a base top-degree form and its second term is
a wedge of closed forms. Therefore $\gamma_\epsilon$ is admissible.
Smooth Stokes for $\varphi_\epsilon$ gives exactly
\begin{align*}
 \int_X\theta_\epsilon\wedge T_\epsilon
 &=\int_X\Omega_\epsilon\wedge T_\epsilon
 =c_\epsilon\int_X(a+\epsilon)^{-1}\mu=n,\\
 D_\mu(\gamma_\epsilon)&=c_\epsilon^{m/n}H_\epsilon.
\end{align*}
Their quotient is $c_\epsilon^{1/n}$. For any smooth Gauduchon test,
the smooth positive representative $\Omega_\epsilon$ and matrix AM--GM
give the matching lower bound $c_\epsilon^{1/n}$. Hence
$V_\mu(\theta_\epsilon)=n/H_\epsilon$.
Since $\theta\leq\theta_\epsilon$, the same tests have unperturbed
quotient at most $c_\epsilon^{1/n}$. Letting $\epsilon\downarrow0$
and using \eqref{eq:lower} proves $V_\mu(\theta)=n$.

The decreasing nef definition with $\theta+\eta h$ agrees as well.
Its infima are at least $\lambda_\mu(\theta)$; for every fixed admissible
$\gamma$ its perturbed quotient converges to the unperturbed quotient.
Taking the infimum in this fixed-test upper bound proves the reverse limit
inequality. No smooth minimizer for the endpoint itself is asserted.

\section{Failure of bigness against all quasi-psh potentials}
Suppose $\theta+\ddc v\geq\delta h$ for a global genuine quasi-psh $v$
and $\delta>0$. Average over the holomorphic transitive fiber action
of $G=\mathrm{SU}(2)$ with Haar probability. The forms $\theta,h$ and
probability $\mu$ are invariant. Since $v\in L^1(\mu)$, Fubini defines
its average as an $L^1$ distribution; differentiation commutes with averaging
and the positive-current inequality is preserved.
The average is $q\circ\operatorname{pr}_1$ almost everywhere, where
$q\in L^1(B)$ is the fiber integral of $v$. Therefore
\[
 \theta+\ddc(q\circ\operatorname{pr}_1)-\delta h\geq0.
\]
Wedge this positive current with the base top form $\omega^m$, whose
pullback is strongly semipositive. All horizontal terms vanish by base dimension, including the
distributional $\ddc(q\circ\operatorname{pr}_1)$. The resulting measure is
\[
 (a-\delta)\mu\geq0.
\]
It would imply $a\geq\delta$ almost everywhere. But $a$ is smooth,
$a(p)=0$, and $a<\delta$ on a nonempty open neighborhood of $p$.
This contradiction excludes every strict current representative.
No singular-current restriction to a fiber or alteration of a canonical
potential was used. This proves nonbigness.

\section{Realization as the endpoint of an actual smooth flow}
The averaging argument also shows the following necessary condition:
if $c\omega+d w+\ddc v$ is positive for a real constant $c$ and a smooth
base function $d$, then wedging its fiber average with $\omega^m$ gives
\begin{equation}\label{eq:vertical}
 d\mu\geq0,\qquad d\geq0\text{ almost everywhere}.
\end{equation}
This remains a distributional argument for every global quasi-psh $v$.
Choose the smooth strictly positive initial Hermitian metric
\[
 \omega_0=(n+1)\omega+(a+2)w.
\]
Its volume is $n(n+1)^m(a+2)\mu$. Fubini--Study has normalized Ricci
forms $n\omega$ on $\mathbb P^m$ and $2w$ on $\mathbb P^1$.
Consequently
\[
 \operatorname{Ric}(\omega_0)=n\omega+2w-\ddc\log(a+2),
\]
and its actual Chern--Ricci background is
\begin{equation}\label{eq:background}
 \alpha_t=(n+1-nt)\omega+(a+2-2t)w+t\ddc\log(a+2).
\end{equation}
For $0\leq t<1$, the smooth potential $-t\log(a+2)$ makes this
background strictly positive: the horizontal coefficient is at least one,
and the vertical coefficient is at least $2(1-t)>0$.
At $t=1$,
\begin{equation}\label{eq:endpoint}
 \alpha_1=\theta+\ddc\log(a+2).
\end{equation}
For every $t>1$, a positive current representative of $\alpha_t$ would,
after this smooth gauge removal and fiber averaging, imply
$a+2-2t\geq0$ almost everywhere by \eqref{eq:vertical}.
Its negative values on an open neighborhood of $p$ make this impossible.
In particular no smooth strictly positive representative exists for $t>1$.
At $t=1$ a smooth strict representative would be big, also impossible.

We now apply the exact maximal-time theorem \cite[Theorem 1.2]{TW}.
On any compact complex manifold with smooth positive Hermitian initial
metric, it gives the unique smooth Chern--Ricci solution until the
supremum of times when the background has a smooth strict potential
representative. It requires no bounded-mass, Gauduchon, pluriclosed or
Guan--Li assumption. The explicit representatives and obstruction above
make that supremum exactly one in our convention.
The source uses $\operatorname{Ric}_{\mathrm{std}}
=-i\partial\bar\partial\log\det g=2\pi\operatorname{Ric}$.
Thus $t=2\pi s$ for source time $s$ and
$\omega(t)=\omega_{\mathrm{std}}(t/(2\pi))$ solves our flow.
The source maximal time is $1/(2\pi)$, hence ours is $\tau=1$.
Its smooth-potential condition is equivalent after multiplying potentials
by $2\pi$. This is an application to the actual flow from $\omega_0$,
not merely a line of backgrounds designated as a flow.

Equation~\eqref{eq:endpoint} identifies the nef limiting background.
Smooth gauge gives exact potential bijections, so it preserves nefness
and bigness. Distributional Stokes against every smooth Gauduchon test
preserves each numerator, including after positive regularization.
The preceding exact volume computation and \eqref{eq:endpoint} give
$V_\mu(\alpha_1)=n>0$, although $\alpha_1$ is not big.
This establishes the numerical counterexample in Theorem~\ref{th:flow}.

\section{Annihilating currents and averaged degeneration}
\begin{lemma}\label{lem:averaged}
For any actual smooth Chern--Ricci flow on $[0,\tau)$, where
$0<\tau<\infty$, let $T$ be a positive $(n-1,n-1)$ current with
$\ddc T=0$, $\int_X\omega_0\wedge T=1$ and
$\int_X(\omega_0-\tau\operatorname{Ric}(\omega_0))\wedge T=0$.
Then
\[
 \int_X\omega(t)\wedge T=1-t/\tau,\qquad 0\leq t<\tau.
\]
\end{lemma}
\begin{proof}
Set $\alpha_t=\omega_0-t\operatorname{Ric}(\omega_0)$.
The flow equation gives the global smooth gauge
\[
 \omega(t)=\alpha_t+\ddc\phi_t,\qquad
 \phi_t=\int_0^t\log\frac{\omega(s)^n}{\omega_0^n}\,ds.
\]
The volume ratio is a global smooth positive function on every compact
time interval before $\tau$. Distributional Stokes against the fixed
$\ddc$-closed current removes the smooth potential. The endpoint pairing
gives $\int_X\operatorname{Ric}(\omega_0)\wedge T=1/\tau$, proving the
identity. No smoothness or $d$-closedness of $T$ is required.
\end{proof}
This identity gives averaged trace degeneration. It does not by itself
give analytic support, pointwise collapse, local convergence or completeness.

The analytic fiber $Z=\{p\}\times\mathbb P^1$ has complex dimension one.
Its integration current has bidegree $(n-1,n-1)$ and is positive closed.
Put $T=\frac12[Z]$. Restrictions here are only of smooth forms:
$\omega_0|_Z=2w$ and $\alpha_1|_Z=0$. Therefore
\[
 \ddc T=0,\qquad \int_X\omega_0\wedge T=1,
 \qquad \int_X\alpha_1\wedge T=0.
\]
Its trace probability is $\sigma_T=\omega_0\wedge T$ and is supported
on the explicitly given analytic fiber. In an $\omega_0$-unitary frame
along $Z$, let $\Pi_{\mathrm{vert}}$ be orthogonal projection onto its
complex tangent line. In the convention
\[
 \beta\wedge T=\frac1n\operatorname{tr}(B_\beta P_T)\,d\sigma_T,
 \qquad\operatorname{tr}P_T=n,
\]
we have $P_T=n\Pi_{\mathrm{vert}}$ for $\sigma_T$-almost every point.
Its rank is one, and $\det P_T=0$. Thus for any smooth positive reference
probability $dV$, the determinant measure
\[
 \Gamma_0(T)=(dV/\omega_0^n)^{1/n}(\det P_T)^{1/n}\sigma_T
\]
is identically zero. Positive fixed-density volume over all smooth
positive tests therefore coexists with this singular annihilator.
There is no claim that rank deficiency in general forces analytic support;
support here follows from the actual definition $T=[Z]/2$.

In fact every normalized annihilator $S$ of this particular endpoint has
trace supported on $Z$ and the same rank-one Radon matrix there. Smooth
gauge Stokes gives $\int_X\theta\wedge S=0$. The measure
$\theta\wedge S$ is nonnegative, so it is identically zero. On every
compact subset $K$ of $X\setminus Z$, one has
$\theta\geq c_K\omega_0$ for some $c_K>0$. Thus its trace
$\sigma_S=\omega_0\wedge S$ vanishes on $K$, and exhaustion gives
$\operatorname{supp}\sigma_S\subset Z$.
In the stated Radon-matrix convention, positive-current coefficient
calculus \cite[III.1.13--1.16, pp.~132--134]{Demailly} makes all coefficient
measures absolutely continuous with respect to $\sigma_S$, giving
$P_S\geq0$ and $\operatorname{tr}P_S=n$ almost everywhere.
Along $Z$ the matrix of $\theta$ is strictly positive on the horizontal
$m$-plane and zero on the vertical line. Its zero pairing with $P_S$
forces every horizontal diagonal entry of $P_S$ to vanish. The inequality
$|(P_S)_{ij}|^2\leq(P_S)_{ii}(P_S)_{jj}$ then removes their rows and
columns. The remaining vertical entry is $n$, so
$P_S=n\Pi_{\mathrm{vert}}$ and $\Gamma_0(S)=0$.
This argument uses the explicit semipositive representative and its kernel;
it is not an inference from determinant vanishing alone.
Lemma~\ref{lem:averaged} specializes to $\int_X\omega(t)\wedge T=1-t$
for the named current. No uniqueness of annihilators is asserted.

\begin{remark}
The non-Hermitian locus of a big background is not assigned to this
nonbig endpoint. We neither refute the local convergence theorems that
assume bigness nor prove weak continuation, completeness near $Z$, or
detailed collapse of the evolving metrics. The base dimension $m\geq2$
is essential for the bounded radial potential and the integrable
$s^{-1}$ determinant majorant. The example is in dimension at least
three and has a genuinely nonclosed initial background. No surface
counterexample or literature-priority claim is made.
\end{remark}

\begin{thebibliography}{9}
\bibitem{Demailly} J.-P. Demailly, \emph{Complex Analytic and Differential
Geometry}, manuscript, 21 June 2012, III.1.13--1.16, pp.~132--134;
\url{https://people.math.harvard.edu/~demarco/Math274/Demailly.pdf}.
\bibitem{TW} V. Tosatti and B. Weinkove, \emph{On the evolution of a Hermitian
metric by its Chern--Ricci form}, arXiv:1201.0312v3 (19 June 2014),
Theorem 1.2 and Sections 3--4; \url{https://arxiv.org/abs/1201.0312}.
\end{thebibliography}
\end{document}
